For every $S$-scheme $X$, we shall denote by $\pi$ or $\pi_X$ the structure morphism $X\to S$.

Given a property $\mathscr P(u)$ for a morphism of $S$-schemes $u:X\to Y$, we shall say that $\mathscr P(u)$ is stable under base change if, whenever $u$ satisfies $\mathscr P(u)$, the same holds for the morphism $u_{(Y')}$, for every $S$-morphism $Y'\to Y$. One says that $\mathscr P(u)$ is local for the topology $\mathcal T$ (cf. Exp. IV, §§4 and 6) if $\mathscr P(u)$ satisfies the following two conditions:
\begin{enumeratea}
\item $\mathscr P(u)$ is stable under base change,
\item whenever there is a family of $S$-morphisms $Y_i\to Y$ covering for the topology $\mathcal T$ and such that each of the morphisms $u_{(Y_i)}$ satisfies $\mathscr P(u)$, then $u$ satisfies $\mathscr P(u)$.
\end{enumeratea}
\begin{proposition}{2.1}\label{VIB.2.1}
Let $\mathscr P(u)$ be a property for a morphism of $S$-schemes, local for the (fpqc) topology. Let $u:G\to H$ be a morphism of $S$-groups. Suppose $G$ flat and universally open over $S$.

Let $W$ be the largest open subset of $H$ above which $u$ satisfies the property $\mathscr P(u)$ and let $V=u^{-1}(W)$. Then $U=\pi_G(V)$ is an open subset of $S$ and $V$ is an open subgroup scheme of $G|_U$.
\end{proposition}
The existence of a largest open subset $W$ of $H$ above which $u$ satisfies $\mathscr P(u)$ follows from the fact that $\mathscr P(u)$ is local for the Zariski topology. Since $\pi_G$ is universally open, $\pi_G(V)$ is an open subset $U$ of $S$. It suffices to show that $V$ is a subgroup scheme of $G|_U$. We may therefore suppose that $U=S$.

Then put $G'=G\times_SV$, $H'=H\times_SV$, $V'=V\times_SV$, $W'=W\times_SV$ and $u'=u_{(V)}$; let $W'_1$ be the largest open subset of $H'$ above which $u'$ satisfies $\mathscr P(u)$; since $V$ is flat and universally open over $S$, the same holds for $H'$ over $H$, and Lemma 2.1.1 below shows that $W'_1=W'$. Then consider the automorphism of $V$-schemes $a$ (resp. $b$) of $G'$ (resp. $H'$), right translation by the inverse of the diagonal section $\delta$ (resp. by the inverse of $u(\delta)$), defined by
\[
a(g,v)=(g\cdot v^{-1},v),\qquad
(\text{resp. }(h,v)=(h\cdot u(v^{-1}),v)),
\]
for every morphism $T\to S$, $g\in G(T)$, $v\in V(T)$ and $h\in H(T)$. It is clear that $u'\circ a=b\circ u'$, which shows that $W'$ is stable under $b$, hence that $V'$ is stable under $a$, so that $V$ is a subgroup scheme of $G$.
% typo?: kept as printed: b is omitted before (h,v) in the displayed formula.
\begin{lemma}{2.1.1}\label{VIB.2.1.1}
Let $\mathscr P(u)$ be a property for an $S$-morphism $u$, local for the (fpqc) topology. Consider a cartesian square of morphisms of $S$-schemes:
\[
\xymatrix{X'\ar[r]^{f'}\ar[d]&Y'\ar[d]^g\\X\ar[r]^f&Y}\ ,
\]
where $g$ is flat and open. Let $W$ (resp. $W'_1$) be the largest open subset of $Y$ (resp. $Y'$) above which $f$ (resp. $f'=f_{(Y')}$) satisfies $\mathscr P(u)$. Then $W'_1=W\times_YY'$.
\end{lemma}
Put $W'=W\times_YY'$; since $\mathscr P(u)$ is stable under base change, one has $W'\subset W'_1$. Since $g$ is open, $W_1=g(W'_1)$ is an open subset of $Y$. Put $V_1=f^{-1}(W_1)$ and $V'_1=V_1\times_{W_1}W'_1$; it is clear that $V'_1=f'^{-1}(W'_1)$. Since $g$ is flat and open, the morphism $W'_1\to W_1=g(W'_1)$ obtained from $g$ is faithfully flat and open, hence covering for the (fpqc) topology (cf. IV 6.3.1 (iv)). Since the morphism $V'_1\to W'_1$ obtained from $f'$ satisfies $\mathscr P(u)$, the same holds for the morphism $V_1\to W_1$ obtained from $f$, hence $W_1\subset W$, and $W'_1\subset g^{-1}(W_1)\subset g^{-1}(W)=W'$; thus $W'=W'_1$.
\begin{remark}{2.1.2}\label{VIB.2.1.2}
A great many properties for a morphism are local for the (fpqc) topology; let us mention those of being flat, (universally) open, (locally) of finite type, of finite presentation, quasi-finite (cf. EGA IV$_2$, 2.5.1, 2.6.1 and 2.7.1), smooth, étale, unramified (EGA IV$_4$, 17.7.3).

The proof of 2.1 in fact uses only base changes by flat morphisms; the proposition will therefore apply to a property satisfying condition b) of 2.0 for the (fpqc) topology, and stable under base changes by flat morphisms (example: that of being quasi-compact and dominant).

Of course, one may state an analogous proposition concerning properties local for a topology $\mathcal T$ finer than the Zariski topology, the condition to be checked on $G$ then being that $\pi_G$ is universally open and covering for the topology $\mathcal T$.

In particular, if $G$ is flat and locally of finite presentation over $S$, one has an analogous statement for properties stable under base changes by flat morphisms locally of finite presentation, and satisfying condition b) of 2.0 for the (fpqc) topology (for example, those of being regular, reduced, Cohen--Macaulay, etc. (EGA IV$_2$, 6.8)).
\end{remark}
\begin{proposition}{2.2}\label{VIB.2.2}
Let $G$ and $H$ be two $S$-groups and $u:G\to H$ a morphism of $S$-groups. Then:\nde{In case (i), the open subset $V$ consists of all the points of $G$ at which $u$ is smooth, resp. étale, resp. of finite presentation and flat, cf. N.D.E. (6). On the other hand, in case (ii), $V$ is the largest open subset contained in the set $E$ of points of $G$ at which $u$ is universally open, but $E$ is not necessarily open, as the following example shows (EGA IV$_3$, 14.1.3 (i)): let $k$ be a field, $H=S=\Spec A$, where $A=k[x]$, and $G$ the $S$-group $\Spec A[y]/(xy)$; then $E$ equals the unit section of $G$, which is not open.}
\begin{enumeratei}
\item Suppose $G$ or $H$ flat over $S$, and $G$ or $H$ locally of finite presentation over $S$, and let $V$ be the largest open subset of $G$ such that the restriction of $u$ to $V$ is flat and locally of finite presentation (resp. smooth, resp. étale). Then $U=\pi_V(V)$ is an open subset of $S$, and $V$ is an open subgroup scheme of $G|_U$.
\item Suppose $G$ or $H$ universally open over $S$, and let $V$ be the largest open subset of $G$ such that the restriction of $u$ to $V$ is universally open. Then $U=\pi_V(V)$ is an open subset of $S$, and $V$ is an open subgroup scheme of $G|_U$.
\end{enumeratei}
\end{proposition}
Let us first prove (i). Let us show that the restriction $\pi_V$ of $\pi_G$ to $V$ is flat and locally of finite presentation:
\begin{enumeratea}
\item If $\pi_G$ is flat (resp. locally of finite presentation), so is $\pi_V$.
\item If $\pi_H$ is flat (resp. locally of finite presentation), so is $\pi_V$, as the composite of the restriction of $u$ to $V$ and $\pi_H$.
\end{enumeratea}
Thus, in the four cases considered in the statement, $\pi_V$ is flat and locally of finite presentation, hence universally open (EGA IV$_2$, 2.4.6). Put $U=\pi_V(V)$; $U$ is therefore an open subset of $S$. It then suffices to show that $V$ is an open subgroup scheme of $G|_U$; we may therefore suppose that $U=S$.

Then put $G'=G\times_SV$, $H'=H\times_SV$, $V'=V\times_SV$ and $u'=u_{(V)}$. Then, since $V$ is flat and locally of finite presentation over $S$, the same holds for $H'$ over $H$. By EGA IV$_4$, 17.7.4, $V'$ is then the largest open subset of $G'$ such that the restriction of $u'$ to $V'$ is flat and locally of finite presentation (resp. smooth, resp. étale). With the automorphisms $a$ and $b$ defined as in the proof of 2.1, it is then clear that $V'$ is stable under $a$, hence that $V$ is a subgroup scheme of $G$.

Let us prove (ii). The restriction $\pi_V$ of $\pi_G$ to $V$ is a universally open morphism, either because this holds for $\pi_G$, or as the composite of the restriction of $u$ to $V$ and $\pi_H$ in the case where $\pi_H$ is universally open. Put $U=\pi_V(V)$; $U$ is then an open subset of $S$. It suffices to show that $V$ is an open subgroup scheme of $G|_U$. We may therefore suppose that $U=S$.

As before put $G'=G\times_SV$, $H'=H\times_SV$, $V'=V\times_SV$ and $u'=u_{(V)}$. Then $\pi_V:V\to S$ is surjective and universally open, and so is $G'\to G$, so that $V'$ is the largest open subset of $G'$ such that the restriction of $u'$ to $V'$ is universally open, by (EGA IV$_3$, 14.3.4 (i) and (ii)). With the automorphisms $a$ and $b$ defined as before, it is then clear that $V'$ is stable under $a$, hence that $V$ is a subgroup scheme of $G$.
\begin{corollary}{2.3}\label{VIB.2.3}
Let $G$ be an $S$-group and $V$ the largest open subset of $G$ flat and locally of finite presentation (resp. smooth, étale, universally open) over $S$. Then $U=\pi_V(V)$ is an open subset of $S$, and $V$ is an open subgroup scheme of $G|_U$.
\end{corollary}
It suffices to apply 2.2 to the case where $H$ is the unit $S$-group and $u$ is the unique morphism of $S$-groups $G\to H$, for then $\pi_H$ is an isomorphism and $\pi_G=\pi_H\circ u$.
\begin{corollary}{2.4}\label{VIB.2.4}
Let $G$ be an $S$-group; if there is a neighborhood $X$ of the unit section having one (or several) of the following properties:
\begin{quote}
$X$ is flat and locally of finite presentation (resp. smooth, étale, universally open) over $S$,
\end{quote}
then there is an open subgroup scheme of $G$ having the same properties.
\end{corollary}
It suffices to apply 2.3, observing that here, with the notation of 2.2, one has $\varepsilon(S)\subset V$, hence $U=S$.
\begin{proposition}{2.5}\label{VIB.2.5}\nde{The statement and proof have been expanded, and in (i) the hypothesis on $H$ has been weakened by replacing ``of finite presentation'' by ``of finite type''.}
Let $u:G\to H$ be a morphism of $S$-groups.
\begin{enumeratei}
\item Suppose that $G$ (resp. $H$) is of finite presentation and flat (resp. of finite type) over $S$ at the points of its unit section. Then the sets
\[
U_{\mathrm{flat}}\supset U_{\mathrm{smooth}}\supset U_{\text{ét.}},
\]
formed by the points $s\in S$ such that $u_s$ is flat (resp. smooth, étale), are open in $S$.

If moreover $G$ (resp. $H$) is locally of finite presentation and flat (resp. locally of finite type) over $S$, then the set $V_{\mathrm{flat}}$ (resp. $V_{\mathrm{smooth}}$, $V_{\text{ét.}}$) of points of $G$ where $u$ is flat (resp. smooth, resp. étale) is the inverse image under $\pi_G$ of $U_{\mathrm{flat}}$ (resp. $U_{\mathrm{smooth}}$, $U_{\text{ét.}}$).
\item Suppose that, for every $s\in S$, $G_s$ is locally of finite type over $\kappa(s)$ (a condition satisfied if $G$ is of finite type over $S$ at the points of the unit section (1.3.1.1)), and that $u$ is locally of finite type (resp. locally of finite presentation) at the points of the unit section of $G$. Then the sets
\[
U_{\mathrm{l.q.f.}}\supset U_{\mathrm{unr.}}
\]
formed by the $s\in S$ such that $u_s$ is locally quasi-finite (resp. unramified) are open in $S$.

If moreover $u$ is locally of finite type (resp. locally of finite presentation), then the set $V_{\mathrm{l.q.f.}}$ (resp. $V_{\mathrm{unr.}}$) of points of $G$ where $u$ is quasi-finite (resp. unramified) is the inverse image under $\pi_G$ of $U_{\mathrm{l.q.f.}}$ (resp. $U_{\mathrm{unr.}}$).
\end{enumeratei}
\end{proposition}
Let us prove (i). First note (1.3.1.1) that, for every $s\in S$, $G_s$ is locally of finite type over $\kappa(s)$. Let $Y$ be the open subset of $H$ formed by the points where $\pi_H$ is of finite type, and let $X_1$ be the open subset of $u^{-1}(Y)$ formed by the points where $\pi_G$ is of finite presentation. By EGA IV$_3$, 11.3.1, the set $X$ of points of $X_1$ where $\pi_G$ is of finite presentation and flat is open in $X_1$, hence in $G$.

Let $u_X:X\to Y$ be the morphism obtained from $u$. Since $Y$ (resp. $X$) contains the unit section of $H$ (resp. $G$), the restriction $\pi_X$ of $\pi_G$ to $X$ is surjective, and the morphism $S\to X$ obtained from $\varepsilon_G$ is an $S$-section of $X$.

Consider the sets $W_{\mathrm{flat}}\supset W_{\mathrm{smooth}}\supset W_{\text{ét.}}$, formed by the points of $X$ where $u_X$ is flat (resp. smooth, resp. étale). It is clear that $W_{\mathrm{smooth}}$ and $W_{\text{ét.}}$ are open in $X$, cf. N.D.E. (6). On the other hand, since $\pi_Y$ is locally of finite type and $\pi_X=\pi_Y\circ u_X$ locally of finite presentation, by EGA IV$_1$, 1.4.3 (v), $u_X$ is locally of finite presentation. Consequently, the flatness locus $W_X$ of $u_X$ is open in $X$ (EGA IV$_3$, 11.3.1).
% typo?: kept as printed: W_X denotes the flatness locus here.

Let $x\in X$ and put $s=\pi_X(x)$; then, by EGA IV$_2$, 11.3.10 (combined with EGA IV$_4$, 17.5.1, resp. 17.6.1), $x$ belongs to $W_{\mathrm{flat}}$ (resp. $W_{\mathrm{smooth}}$, $W_{\text{ét.}}$) if and only if $u_s$ is flat (resp. smooth, étale) at the point $x$, or, equivalently by 1.3, if and only if $u_s$ is flat (resp. smooth, étale).
% typo?: kept as printed: EGA IV_2, 11.3.10.

Consequently, for ``$\flat=\text{flat, smooth or étale}$'', one has $U_\flat=\varepsilon_G^{-1}(W_\flat)$ and $W_\flat=\pi_X^{-1}(U_\flat)$. Since $W_\flat$ is open in $X$, hence in $G$, the first equality shows that $U_\flat$ is open in $S$.

The second assertion of (i) follows from what has just been seen, for then $Y=H$, $X=G$, and $V_\flat=W_\flat=\pi_G^{-1}(U_\flat)$.

Let us prove assertion (ii). Let $V_{\mathrm{l.f.t.}}\supset V_{\mathrm{l.q.f.}}$ be the open subsets of $G$ formed by the points at which $u$ is of finite type (resp. quasi-finite). Put $X=V_{\mathrm{l.f.t.}}$ and denote by $u_X$ the restriction of $u$ to $X$. By hypothesis, $X$ contains $\varepsilon(S)$. Let $x\in X$ and put $s=\pi_X(x)$.

If $u$ is quasi-finite at $x$, then, by base change, $u_s$ is quasi-finite at $x$, and hence, since $G_s$ is supposed locally of finite type, $u_s$ is locally quasi-finite, by 1.3.1.

Conversely, if $u_s$ is locally quasi-finite, then $u_X^{-1}(u_X(x))\subset u_s^{-1}(u_s(x))$ is a finite set. Since $u_X:X\to H$ is locally of finite type, it follows from Chevalley's semicontinuity theorem (EGA IV$_3$, 13.1.3) that there is an open neighborhood $W$ of $x$ in $X$ such that the fiber $u_X^{-1}(u_X(x'))$ is discrete for every $x'\in W$. Thus, by EGA II, 6.2.2 (and EGA III$_2$, ErrIII 20), $u_X$ is quasi-finite at $x$.

Consequently, one has $U_{\mathrm{l.q.f.}}=\varepsilon_G^{-1}(V_{\mathrm{l.q.f.}})$ and $V_{\mathrm{l.q.f.}}=\pi_X^{-1}(U_{\mathrm{l.q.f.}})$, and the first equality shows that $U_{\mathrm{l.q.f.}}$ is open in $S$.

If moreover $G$ is locally of finite type over $S$, then $X=G$, and hence $V_{\mathrm{l.q.f.}}$ is the inverse image under $\pi_G$ of $U_{\mathrm{l.q.f.}}$. This proves the first half of (ii).
% typo?: kept as printed: G locally of finite type here rather than u.

Now consider the open subsets $V_{\mathrm{l.f.p.}}\supset V_{\mathrm{unr.}}$, formed by the points where $u$ is of finite presentation (resp. unramified), and suppose that $Z=V_{\mathrm{l.f.p.}}$ contains the unit section of $G$. Let $z\in Z$; put $s=\pi_Z(z)$ and $h=u_Z(z)$.

If $u$ is unramified at $z$, then, by base change, $u_s$ is unramified at $z$, and hence, since $G_s$ is supposed locally of finite type, $u_s$ is unramified, by 1.3.1.

Conversely, if $u_s$ is unramified at the point $z$, then the fiber $u_s^{-1}(h)=u^{-1}(h)$ is unramified over $\kappa(h)$ at the point $z$. Since $Z$ is an open subset of $G$, the fiber $u_Z^{-1}(h)=u_s^{-1}(h)\cap Z$ is unramified over $\kappa(h)$ at the point $z$. Since $u_Z$ is locally of finite presentation, this implies, by EGA IV$_4$, 17.4.1, that $u_Z$ is unramified at the point $z$.

Thus one has $U_{\mathrm{unr.}}=\varepsilon_G^{-1}(V_{\mathrm{unr.}})$ and $V_{\mathrm{unr.}}=\pi_Z^{-1}(U_{\mathrm{unr.}})$, and the first equality shows that $U_{\mathrm{unr.}}$ is open in $S$.

If moreover $G$ is locally of finite presentation over $S$, then $Z=G$, and hence $V_{\mathrm{unr.}}$ is the inverse image under $\pi_G$ of $U_{\mathrm{unr.}}$. This completes the proof of assertion (ii).
% typo?: kept as printed: G locally of finite presentation here rather than u.
\begin{corollary}{2.6}\label{VIB.2.6}
Let $u:G\to H$ be a morphism of $S$-groups which is a radicial morphism (as is the case if $u$ is a monomorphism (EGA I, 3.5.4)). Suppose $G$ (resp. $H$) locally of finite presentation and flat (resp. locally of finite type) over $S$. Then the set $U$ of $s\in S$ such that $u_s$ is an open immersion is open in $S$, and the restriction of $u$ above $U$ is an open immersion.
\end{corollary}
By 2.5 (i), the set $U'$ of points $s\in S$ such that $u_s$ is étale is open in $S$. Since $u$ is radicial, so is $u_s$, for every $s\in S$, hence, by EGA IV$_4$, 17.9.1, one has $U=U'$, which shows that $U$ is open. Finally, by 2.5 (i), the restriction of $u$ above $U$ is étale; since $u$ is radicial, this restriction is an open immersion (EGA IV$_4$, 17.9.1).
\begin{proposition}{2.7}\label{VIB.2.7}
Let $G$ be an $S$-group. The following conditions are equivalent:
\begin{enumeratei}
\item $G$ is unramified over $S$ at the points of the unit section,
\item the unit section is an open immersion,
\item $G$ is of finite presentation over $S$ at the points of the unit section, and, for every $s\in S$, $G_s$ is unramified over $\kappa(s)$.
\end{enumeratei}
If moreover $G$ is supposed locally of finite presentation over $S$, each of the preceding three conditions is equivalent to the following:
\begin{enumeratei}\setcounter{enumi}{3}
\item $G$ is unramified over $S$.
\end{enumeratei}
\end{proposition}
The equivalence of assertions (i) and (ii) follows from the more general Lemma 2.7.1 below. Observe (1.3.1.1) that either of conditions (i) or (iii) implies that, for every $s\in S$, $G_s$ is locally of finite type over $\kappa(s)$. Then (EGA IV$_4$, 17.4.1), assertion (i) is equivalent to the fact that, for every $s\in S$, $G_s$ is unramified over $\kappa(s)$ at the point $e_s$, the identity of $G_s$, or again (1.3.1), to the fact that $G_s$ is unramified over $\kappa(s)$, so assertions (i) and (iii) are equivalent. Finally, if $G$ is locally of finite presentation over $S$, assertions (iii) and (iv) are equivalent (EGA IV$_4$, 17.4.1).
\begin{lemma}{2.7.1}\label{VIB.2.7.1}
Let $G$ be an $S$-scheme endowed with a section $\varepsilon$. For $G$ to be unramified over $S$ at the points of this section, it is necessary and sufficient that $\varepsilon$ be an open immersion.
\end{lemma}
The condition is necessary, by EGA IV$_4$, 17.4.1 a) $\Rightarrow$ b$''$). Conversely, if $\varepsilon$ is an open immersion, the restriction to $\varepsilon(S)$ of the structure morphism $G\to S$ is an isomorphism, so $G$ is unramified over $S$ at the points of $\varepsilon(S)$.
\begin{corollary}{2.8}\label{VIB.2.8}
Let $u:G\to H$ be a morphism of $S$-groups. Suppose that, for every $s\in S$, $G_s$ is locally of finite type over $\kappa(s)$.\nde{The presentation has been modified, separating the assertions concerning the ``quasi-finite'' case from those concerning the ``unramified'' case.}
\begin{enumeratea}
\item If $u$ is locally of finite type, the following conditions are equivalent:
\begin{enumeratei}
\item $u$ is locally quasi-finite,
\item for every $s\in S$, $u_s:G_s\to H_s$ is locally quasi-finite,
\item $\Ker u$ is locally quasi-finite over $S$,
\item the fibers of $\Ker u$ are discrete.
\end{enumeratei}
\item If $u$ is locally of finite presentation, the following conditions are equivalent:
\begin{enumeratei}\setcounter{enumii}{4}
\item $u$ is unramified,
\item for every $s\in S$, $u_s:G_s\to H_s$ is unramified,
\item $\Ker u$ is unramified,
\item the unit section $S\to\Ker u$ is an open immersion.
\end{enumeratei}
\end{enumeratea}
\end{corollary}
The equivalences (i) $\Leftrightarrow$ (ii) and (v) $\Leftrightarrow$ (vi) follow from 2.5 (ii), and the reminder 2.8.1 below shows that (i) $\Rightarrow$ (iii) and (v) $\Rightarrow$ (vii).

For every $s\in S$, denote by $e_s$ the identity element of $H_s$. Then (iii) (resp. (vii)) implies that, for every $s\in S$,
\[
(\Ker u)_s=\Ker u_s=u_s^{-1}(e_s)
\]
is locally quasi-finite (resp. unramified) over $\kappa(s)=\kappa(e_s)$, hence that $u_s$ is quasi-finite (resp. unramified) at the identity point of $G_s$, hence, by 1.3, that $u_s$ is locally quasi-finite (resp. unramified). Thus (iii) $\Rightarrow$ (ii) and (vii) $\Rightarrow$ (vi). Finally, (ii) $\Leftrightarrow$ (iv) by 1.4.1, and (vii) $\Leftrightarrow$ (viii) by 2.7.
\begin{enoncerm}{Reminder}{2.8.1}\label{VIB.2.8.1}
Recall (cf. I 2.3.6.1) that, given a morphism $u:G\to H$ of $S$-groups, one calls the kernel of $u$, and denotes by $\Ker u$, the subgroup functor of $G$ defined by putting, for every morphism $f:T\to S$,
\[
(\Ker u)(T)=\{a\in G(T)\mid u\circ a=\varepsilon_H\circ f\}.
\]
By loc. cit. (or EGA I, 4.4.1), this functor is representable by the $S$-group $G\times_HS=u^{-1}(\varepsilon_H(S))$, denoted simply by $\Ker u$. In particular, the structure morphism $\Ker u\to S$ is obtained from $u$ by base change.
\end{enoncerm}
\begin{lemma}{2.9}\label{VIB.2.9}
Let $\pi:X\to S$ be a morphism admitting an $S$-section $\varepsilon$.
\begin{enumeratei}
\item If $\pi$ is injective, it is integral.\begingroup\renewcommand{\thefootnote}{*}\addtocounter{footnote}{-1}\footnote{This is also a special case of EGA IV$_4$, 18.12.11, for $\pi$ is evidently a universal homeomorphism.}\endgroup
\item If $\pi$ is locally of finite type, and if, for every $s\in S$, $\pi_s$ is an isomorphism, then $\pi$ is an isomorphism.\begingroup\renewcommand{\thefootnote}{**}\addtocounter{footnote}{-1}\footnote{This is also an immediate consequence of EGA IV$_4$, 18.12.6.}\endgroup
\end{enumeratei}
\end{lemma}
First observe that, by Lemma 2.9.1 below, $\pi$, being a homeomorphism, is an affine morphism.

If $\pi$ is injective, $\varepsilon$ is surjective. Since $\varepsilon$ is a surjective immersion, $\varepsilon(S)$ is isomorphic to the closed subscheme of $X$ defined by a nilideal $\mathcal I$ of $\OX$. Since $\varepsilon$ is an $S$-section of the morphism $\pi$, one has a direct sum decomposition $\OX=\OS\oplus\mathcal I$ of $\OS$-modules. Since $\mathcal I$ is a nilideal of $\OX$, $\mathcal I$ is evidently integral over $\OS$, so $\OX$ is integral over $\OS$, and $\pi$ is integral.

Now suppose $\pi$ locally of finite type. Since $\pi\circ\varepsilon=\mathrm{id}_S$, $\varepsilon$ is locally of finite presentation (cf. EGA IV$_1$, 1.4.3 (v)), so $\mathcal I$ is a finitely generated ideal of $\OX$ (EGA IV$_1$, 1.4.1). For every $s\in S$, one has $\OO_{X_s}=\kappa(s)\oplus\mathcal I\otimes_{\OS}\kappa(s)$. By hypothesis, $\varepsilon_s$ is an isomorphism, so $\mathcal I\otimes_{\OS}\kappa(s)=0$ for every $s\in S$, hence a fortiori $\mathcal I\otimes_{\OX}\kappa(x)=0$ for every $x\in X$, which implies, by Nakayama's lemma, that $\mathcal I=0$, hence that $\pi$ is an isomorphism.
\begin{lemma}{2.9.1}\label{VIB.2.9.1}
Let $f:X\to Y$ be a morphism of schemes which is a homeomorphism; then $f$ is an affine morphism.\begingroup\renewcommand{\thefootnote}{***}\addtocounter{footnote}{-1}\footnote{Cf. EGA IV$_4$, 18.12.7.1 for a slightly more general result, proved by the same argument.}\endgroup
\end{lemma}
It suffices to show that every point $y\in Y$ has an open neighborhood $W$ such that the restriction of $f$ above $W$ is an affine morphism. Then let $y\in Y$, and let $V$ be an affine open neighborhood of $y$ in $Y$. Let $V'=f^{-1}(V)$. Then $V'$ is an open neighborhood of $x=f^{-1}(y)$ in $X$. There is an affine open neighborhood $W'$ of $x$ in $X$ contained in $V'$. Then put $W=f(W')$. Then $W$ is an open neighborhood of $y$ in $Y$ contained in the affine scheme $V$, so $W$ is separated. Since $W'$ is an affine scheme, the restriction of $f$ above $W$ is then an affine morphism (EGA II, 1.2.3).
\begin{corollary}{2.10}\label{VIB.2.10}
Let $G$ be an $S$-group locally of finite type. Suppose that, for every $s\in S$, $G_s$ is the unit $\kappa(s)$-group; then $G$ is the unit $S$-group.
\end{corollary}
% typo: quel soit -> quel que soit.
More generally:
\begin{corollary}{2.11}\label{VIB.2.11}
Let $f:X\to Y$ be an $S$-morphism locally of finite type. For $f$ to be a monomorphism, it is necessary and sufficient that, for every $s\in S$, $f_s$ be a monomorphism.
\end{corollary}
It is clear that the condition is necessary; let us show that it is sufficient. If, for every $s\in S$, $f_s$ is a monomorphism, a fortiori for every $y\in Y$, $f_y$ is a monomorphism; we may therefore suppose that $Y=S$.

By EGA I, 5.3.8, to show that $f$ is a monomorphism, it suffices to show that $\Delta_f:X\to X\times_SX$ is an isomorphism, or, equivalently, that the first projection $p:X\times_SX\to X$ is an isomorphism. But if $f_s$ is a monomorphism, it similarly follows from EGA I, 5.3.8 that the first projection $X_s\times_{\kappa(s)}X_s\to X_s$ (which is identified with $p_s$) is an isomorphism. Now $p$ has the $S$-section $\Delta_f$, so Lemma 2.9 asserts that if, for every $s\in S$, $p_s$ is an isomorphism, then so is $p$.
\begin{corollary}{2.12}\label{VIB.2.12}
Let $G$ be an $S$-scheme having an $S$-section $\varepsilon$ and such that the structure morphism $\pi:G\to S$ is closed. Let $s\in S$ be such that $\pi$ is of finite presentation at the point $\varepsilon(s)$, and $\varepsilon_s:\Spec(\kappa(s))\to G_s$ is an isomorphism (or, equivalently, $\pi_s:G_s\to\Spec\kappa(s)$ is an isomorphism). Then there is an open neighborhood $U$ of $s$ in $S$ such that $\varepsilon|_U:U\to G\times_SU$ is an isomorphism.
\end{corollary}
Let $V$ be the set of points of $G$ where $\pi$ is unramified; one knows that $V$ is open (cf. N.D.E. (6)) and contains $\varepsilon(s)$. Thus $W=\varepsilon^{-1}(V)$ is an open subset of $S$ containing $s$, and such that, for every $t\in W$, $\pi$ is unramified at $\varepsilon(t)$. Since $\pi$ is closed, so is $\pi|_W$, so we may suppose that $W=S$.

Then, by 2.7.1, $X=G-\varepsilon(S)$ is a closed subset of $G$ not meeting $G_s$; thus, since $\pi$ is closed, $F=\pi(X)$ is a closed subset of $S$ not meeting $s$; put $U=S-F$; then $U$ is an open subset of $S$ such that $\varepsilon|_U$ is an isomorphism of $U$ onto $G|_U$.

\sgasection{3}{Identity component of a group locally of finite presentation}
\numpara{3.0}Given a subset $A$ (resp. $B$) of an $S$-scheme $X$ (resp. $Y$), by abuse of notation, $A\times_SB$ will denote the subset of $X\times_SY$ formed by the points whose first projection belongs to $A$ and second to $B$.
